% Causal DAGs as world specifications. Numbers come from E64 macros
% (scripts/make_paper_assets.py -> numbers.tex); the table is tables/causal_assumptions.tex.
% Input from main.tex with % Causal DAGs as world specifications. Numbers come from E64 macros
% (scripts/make_paper_assets.py -> numbers.tex); the table is tables/causal_assumptions.tex.
% Input from main.tex with % Causal DAGs as world specifications. Numbers come from E64 macros
% (scripts/make_paper_assets.py -> numbers.tex); the table is tables/causal_assumptions.tex.
% Input from main.tex with % Causal DAGs as world specifications. Numbers come from E64 macros
% (scripts/make_paper_assets.py -> numbers.tex); the table is tables/causal_assumptions.tex.
% Input from main.tex with \input{sections/dag_causal}.

\section{Causal DAGs as World Specifications}\label{sec:dag}
A world model and a causal diagram are the same object viewed from two sides: a
directed acyclic graph (DAG) names variables and asserts which directly cause
which; an \openworld{} \texttt{World} carries a symbolic state ontology and a
transition that recomputes that state. \openworld{} closes the gap by accepting
a DAG as a first-class input and compiling it into a verified world, so that the
causal assumptions a domain expert is willing to commit to in a diagram become
executable, inspectable, and---crucially---falsifiable dynamics.

\subsection{A DAG as a perception modality}
DAGs enter through the same perception boundary (\S\ref{sec:language}) that
admits text, audio, and images: a \texttt{graph} modality whose
\texttt{DAGPerceptor} ingests the \texttt{dagitty}/Graphviz \texttt{.dot} text
those tools already export. A single option selects what the graph resolves
into---the three ways a world model consumes a causal structure. In
\emph{graph} mode it commits a normalized object (nodes with their
exposure/outcome/latent tags, plus edges) to state. In \emph{schema} mode it
emits the observed (non-latent) variables to extract from a record downstream:
the DAG \emph{declares what to perceive}. In \emph{world} mode it compiles the
DAG to a runnable structural causal model (SCM). The perceptor is deterministic
and code-free, so it round-trips losslessly through a portable spec like the
other declarative perceptors.

\subsection{Compiling a DAG to a verified SCM}
Compilation emits a single pure \texttt{transition(state, action)} that
recomputes each node from its parents in topological order, with an exogenous
term $u_v$ per node $v$; the default per-node form is a transparent
linear threshold over the parents, and edge weights/signs are overridable. The
emitted code passes the same gates as any synthesized dynamics
(\S\ref{sec:language})---it parses, smoke-runs sandboxed on every action, and
respects declared invariants---so the causal model is an auditable artifact, not
a black box. Interventions are simply the world's actions: a \texttt{do\_$v$}
action overrides node $v$'s structural equation and the override propagates to
its descendants within the same topological pass, realizing Pearl's do-operator
\cite{pearl2009causality}. Because the dynamics are deterministic given the
exogenous terms, a unit-level counterfactual is recovered exactly: hold a unit's
$u_v$ fixed, swap the action, and re-roll (abduction--action--prediction). The
result is a training-free interventional and counterfactual sandbox over an
\emph{assumed} causal structure.

\subsection{Testing the assumptions, not just running them}
A diagram is a set of hypotheses, and \openworld{} lets a corpus of observed
transitions adjudicate among them on the factored many-worlds store of
\S\ref{sec:active-induction}. Each uncertain edge becomes a parameter whose
candidate values are the rival assumptions---an edge's presence, or the sign of
an effect---and each node a \texttt{Mechanism} predicting a discretized
(cut-point) outcome from its parents under those parameters. Observing
transitions reweights only the small factors the firing mechanisms touch, so the
version space over causal models is pruned without ever enumerating the joint.

On a synthetic cohort of \DagCleanCohort{} units drawn from a known structure
(E64; shaped after the perinatal DAG of \S\ref{sec:dag-app}), exact pruning
collapses \DagCandidateModels{} candidate causal models to \DagSurvivorsClean{},
recovering all \DagIdentifiedEdges{} identified edges
(income~$\to$~stress, stress~$\to$~preterm, preterm~$\to$~outcome) at posterior
$1.0$ (Table~\ref{tab:causal-assumptions}). The \DagSurvivorsClean{} survivors
differ only in whether a \emph{direct} stress~$\to$~outcome edge is present:
with stress fully determining preterm in this cohort, that edge is
\emph{non-identifiable}, and the store reports it as an even posterior rather
than committing---the mediator-versus-direct ambiguity surfaced and measured,
not hidden, in the same spirit as the framework's other boundary results. Hard
version-space pruning is brittle to measurement noise: under \DagNoiseFlipPct\%{}
outcome noise the true model is contradicted and \DagHardSurvivorsNoisy{}
candidates survive. An additive soft-Bayesian scorer over the \emph{same}
mechanisms, weighting each observation by a flip-noise likelihood, instead holds
the identified edges at posterior $\geq\DagSoftMinIdentified{}$ while still
leaving the unidentifiable edge near-even---motivating a likelihood-weighted
semiring as the natural generalization of the store's update.

\begin{table}[htb]
\centering
\caption{\textbf{E64: testing causal assumptions on the many-worlds store.} Each
questioned edge of the perinatal DAG is a parameter; exact version-space pruning
on clean data pins the identified edges to their true value and leaves the
non-identifiable direct edge even, while a soft-Bayesian scorer over the same
mechanisms recovers the identified edges under \DagNoiseFlipPct\%{} outcome noise.}
\label{tab:causal-assumptions}
\input{tables/causal_assumptions}
\end{table}

\subsection{Application: social determinants of health}\label{sec:dag-app}
The motivating use case is individual health-risk counterfactuals over social
determinants of health (SDOH). A person's record is perceived into a symbolic
SDOH state; structural factors (income, housing, food security, social support,
neighborhood) and a place-based population prior drive a mediating psychosocial
layer (perceived stress, pregnancy-related anxiety), which in turn drives a
physiological load (a neuroendocrine/cortisol and an inflammatory term) and
finally an individualized outcome risk. The chain encodes a documented
mechanism---psychosocial stress in pregnancy raising preterm-birth risk through
neuroendocrine and inflammatory pathways, with social support protective and
effects strongest in early gestation \cite{shapiro2013psp}---over a social
gradient in which structural disadvantage, not individual behavior alone, sets
the baseline \cite{marmot2018}. Interventions are the world's actions, so
counterfactual care plans are rollouts under different action sequences: the
framework compares a structural plan, a proximal stress/anxiety-reduction plan,
and their combination by resulting risk and by risk reduction per unit of effort,
and emits a recommendation. Two worked worlds ship as examples (a generic SDOH
world and a perinatal preterm-birth world); both reproduce the qualitative
literature---structural action dominates behavior-only plans, while proximal
stress care matters most early in gestation.

The scope is deliberately bounded. \openworld{} supplies the verified, auditable
counterfactual layer; it consumes a population prior (e.g.\ a geospatial
foundation-model embedding for a locale) as an exogenous input rather than
learning one. And it tests the \emph{implications and internal consistency} of a
causal assumption set---interventional and counterfactual predictions, edge-sign
and functional-form sensitivity, and the identifiability of a target effect---not
the truth of the diagram, which still requires real data and an identification
argument. What it removes is the opacity: every edge is a line of code, every
recommendation a rollout one can replay, and every assumption a hypothesis the
data is allowed to refute.
.

\section{Causal DAGs as World Specifications}\label{sec:dag}
A world model and a causal diagram are the same object viewed from two sides: a
directed acyclic graph (DAG) names variables and asserts which directly cause
which; an \openworld{} \texttt{World} carries a symbolic state ontology and a
transition that recomputes that state. \openworld{} closes the gap by accepting
a DAG as a first-class input and compiling it into a verified world, so that the
causal assumptions a domain expert is willing to commit to in a diagram become
executable, inspectable, and---crucially---falsifiable dynamics.

\subsection{A DAG as a perception modality}
DAGs enter through the same perception boundary (\S\ref{sec:language}) that
admits text, audio, and images: a \texttt{graph} modality whose
\texttt{DAGPerceptor} ingests the \texttt{dagitty}/Graphviz \texttt{.dot} text
those tools already export. A single option selects what the graph resolves
into---the three ways a world model consumes a causal structure. In
\emph{graph} mode it commits a normalized object (nodes with their
exposure/outcome/latent tags, plus edges) to state. In \emph{schema} mode it
emits the observed (non-latent) variables to extract from a record downstream:
the DAG \emph{declares what to perceive}. In \emph{world} mode it compiles the
DAG to a runnable structural causal model (SCM). The perceptor is deterministic
and code-free, so it round-trips losslessly through a portable spec like the
other declarative perceptors.

\subsection{Compiling a DAG to a verified SCM}
Compilation emits a single pure \texttt{transition(state, action)} that
recomputes each node from its parents in topological order, with an exogenous
term $u_v$ per node $v$; the default per-node form is a transparent
linear threshold over the parents, and edge weights/signs are overridable. The
emitted code passes the same gates as any synthesized dynamics
(\S\ref{sec:language})---it parses, smoke-runs sandboxed on every action, and
respects declared invariants---so the causal model is an auditable artifact, not
a black box. Interventions are simply the world's actions: a \texttt{do\_$v$}
action overrides node $v$'s structural equation and the override propagates to
its descendants within the same topological pass, realizing Pearl's do-operator
\cite{pearl2009causality}. Because the dynamics are deterministic given the
exogenous terms, a unit-level counterfactual is recovered exactly: hold a unit's
$u_v$ fixed, swap the action, and re-roll (abduction--action--prediction). The
result is a training-free interventional and counterfactual sandbox over an
\emph{assumed} causal structure.

\subsection{Testing the assumptions, not just running them}
A diagram is a set of hypotheses, and \openworld{} lets a corpus of observed
transitions adjudicate among them on the factored many-worlds store of
\S\ref{sec:active-induction}. Each uncertain edge becomes a parameter whose
candidate values are the rival assumptions---an edge's presence, or the sign of
an effect---and each node a \texttt{Mechanism} predicting a discretized
(cut-point) outcome from its parents under those parameters. Observing
transitions reweights only the small factors the firing mechanisms touch, so the
version space over causal models is pruned without ever enumerating the joint.

On a synthetic cohort of \DagCleanCohort{} units drawn from a known structure
(E64; shaped after the perinatal DAG of \S\ref{sec:dag-app}), exact pruning
collapses \DagCandidateModels{} candidate causal models to \DagSurvivorsClean{},
recovering all \DagIdentifiedEdges{} identified edges
(income~$\to$~stress, stress~$\to$~preterm, preterm~$\to$~outcome) at posterior
$1.0$ (Table~\ref{tab:causal-assumptions}). The \DagSurvivorsClean{} survivors
differ only in whether a \emph{direct} stress~$\to$~outcome edge is present:
with stress fully determining preterm in this cohort, that edge is
\emph{non-identifiable}, and the store reports it as an even posterior rather
than committing---the mediator-versus-direct ambiguity surfaced and measured,
not hidden, in the same spirit as the framework's other boundary results. Hard
version-space pruning is brittle to measurement noise: under \DagNoiseFlipPct\%{}
outcome noise the true model is contradicted and \DagHardSurvivorsNoisy{}
candidates survive. An additive soft-Bayesian scorer over the \emph{same}
mechanisms, weighting each observation by a flip-noise likelihood, instead holds
the identified edges at posterior $\geq\DagSoftMinIdentified{}$ while still
leaving the unidentifiable edge near-even---motivating a likelihood-weighted
semiring as the natural generalization of the store's update.

\begin{table}[htb]
\centering
\caption{\textbf{E64: testing causal assumptions on the many-worlds store.} Each
questioned edge of the perinatal DAG is a parameter; exact version-space pruning
on clean data pins the identified edges to their true value and leaves the
non-identifiable direct edge even, while a soft-Bayesian scorer over the same
mechanisms recovers the identified edges under \DagNoiseFlipPct\%{} outcome noise.}
\label{tab:causal-assumptions}
\begin{tabular}{l l c c c}
\toprule
Edge & Candidates & Exact post.\ & Soft post.\ & Identified \\
& & (clean) & (noisy) & \\
\midrule
Income $\to$ Stress & causal/null & 1.00 & 1.00 & \checkmark \\
Stress $\to$ Preterm & causal/null & 1.00 & 1.00 & \checkmark \\
Preterm $\to$ Bayley & causal/null & 1.00 & 1.00 & \checkmark \\
Stress $\to$ Bayley (direct) & yes/no & 0.50 & 0.50 & -- \\
\bottomrule
\end{tabular}

\end{table}

\subsection{Application: social determinants of health}\label{sec:dag-app}
The motivating use case is individual health-risk counterfactuals over social
determinants of health (SDOH). A person's record is perceived into a symbolic
SDOH state; structural factors (income, housing, food security, social support,
neighborhood) and a place-based population prior drive a mediating psychosocial
layer (perceived stress, pregnancy-related anxiety), which in turn drives a
physiological load (a neuroendocrine/cortisol and an inflammatory term) and
finally an individualized outcome risk. The chain encodes a documented
mechanism---psychosocial stress in pregnancy raising preterm-birth risk through
neuroendocrine and inflammatory pathways, with social support protective and
effects strongest in early gestation \cite{shapiro2013psp}---over a social
gradient in which structural disadvantage, not individual behavior alone, sets
the baseline \cite{marmot2018}. Interventions are the world's actions, so
counterfactual care plans are rollouts under different action sequences: the
framework compares a structural plan, a proximal stress/anxiety-reduction plan,
and their combination by resulting risk and by risk reduction per unit of effort,
and emits a recommendation. Two worked worlds ship as examples (a generic SDOH
world and a perinatal preterm-birth world); both reproduce the qualitative
literature---structural action dominates behavior-only plans, while proximal
stress care matters most early in gestation.

The scope is deliberately bounded. \openworld{} supplies the verified, auditable
counterfactual layer; it consumes a population prior (e.g.\ a geospatial
foundation-model embedding for a locale) as an exogenous input rather than
learning one. And it tests the \emph{implications and internal consistency} of a
causal assumption set---interventional and counterfactual predictions, edge-sign
and functional-form sensitivity, and the identifiability of a target effect---not
the truth of the diagram, which still requires real data and an identification
argument. What it removes is the opacity: every edge is a line of code, every
recommendation a rollout one can replay, and every assumption a hypothesis the
data is allowed to refute.
.

\section{Causal DAGs as World Specifications}\label{sec:dag}
A world model and a causal diagram are the same object viewed from two sides: a
directed acyclic graph (DAG) names variables and asserts which directly cause
which; an \openworld{} \texttt{World} carries a symbolic state ontology and a
transition that recomputes that state. \openworld{} closes the gap by accepting
a DAG as a first-class input and compiling it into a verified world, so that the
causal assumptions a domain expert is willing to commit to in a diagram become
executable, inspectable, and---crucially---falsifiable dynamics.

\subsection{A DAG as a perception modality}
DAGs enter through the same perception boundary (\S\ref{sec:language}) that
admits text, audio, and images: a \texttt{graph} modality whose
\texttt{DAGPerceptor} ingests the \texttt{dagitty}/Graphviz \texttt{.dot} text
those tools already export. A single option selects what the graph resolves
into---the three ways a world model consumes a causal structure. In
\emph{graph} mode it commits a normalized object (nodes with their
exposure/outcome/latent tags, plus edges) to state. In \emph{schema} mode it
emits the observed (non-latent) variables to extract from a record downstream:
the DAG \emph{declares what to perceive}. In \emph{world} mode it compiles the
DAG to a runnable structural causal model (SCM). The perceptor is deterministic
and code-free, so it round-trips losslessly through a portable spec like the
other declarative perceptors.

\subsection{Compiling a DAG to a verified SCM}
Compilation emits a single pure \texttt{transition(state, action)} that
recomputes each node from its parents in topological order, with an exogenous
term $u_v$ per node $v$; the default per-node form is a transparent
linear threshold over the parents, and edge weights/signs are overridable. The
emitted code passes the same gates as any synthesized dynamics
(\S\ref{sec:language})---it parses, smoke-runs sandboxed on every action, and
respects declared invariants---so the causal model is an auditable artifact, not
a black box. Interventions are simply the world's actions: a \texttt{do\_$v$}
action overrides node $v$'s structural equation and the override propagates to
its descendants within the same topological pass, realizing Pearl's do-operator
\cite{pearl2009causality}. Because the dynamics are deterministic given the
exogenous terms, a unit-level counterfactual is recovered exactly: hold a unit's
$u_v$ fixed, swap the action, and re-roll (abduction--action--prediction). The
result is a training-free interventional and counterfactual sandbox over an
\emph{assumed} causal structure.

\subsection{Testing the assumptions, not just running them}
A diagram is a set of hypotheses, and \openworld{} lets a corpus of observed
transitions adjudicate among them on the factored many-worlds store of
\S\ref{sec:active-induction}. Each uncertain edge becomes a parameter whose
candidate values are the rival assumptions---an edge's presence, or the sign of
an effect---and each node a \texttt{Mechanism} predicting a discretized
(cut-point) outcome from its parents under those parameters. Observing
transitions reweights only the small factors the firing mechanisms touch, so the
version space over causal models is pruned without ever enumerating the joint.

On a synthetic cohort of \DagCleanCohort{} units drawn from a known structure
(E64; shaped after the perinatal DAG of \S\ref{sec:dag-app}), exact pruning
collapses \DagCandidateModels{} candidate causal models to \DagSurvivorsClean{},
recovering all \DagIdentifiedEdges{} identified edges
(income~$\to$~stress, stress~$\to$~preterm, preterm~$\to$~outcome) at posterior
$1.0$ (Table~\ref{tab:causal-assumptions}). The \DagSurvivorsClean{} survivors
differ only in whether a \emph{direct} stress~$\to$~outcome edge is present:
with stress fully determining preterm in this cohort, that edge is
\emph{non-identifiable}, and the store reports it as an even posterior rather
than committing---the mediator-versus-direct ambiguity surfaced and measured,
not hidden, in the same spirit as the framework's other boundary results. Hard
version-space pruning is brittle to measurement noise: under \DagNoiseFlipPct\%{}
outcome noise the true model is contradicted and \DagHardSurvivorsNoisy{}
candidates survive. An additive soft-Bayesian scorer over the \emph{same}
mechanisms, weighting each observation by a flip-noise likelihood, instead holds
the identified edges at posterior $\geq\DagSoftMinIdentified{}$ while still
leaving the unidentifiable edge near-even---motivating a likelihood-weighted
semiring as the natural generalization of the store's update.

\begin{table}[htb]
\centering
\caption{\textbf{E64: testing causal assumptions on the many-worlds store.} Each
questioned edge of the perinatal DAG is a parameter; exact version-space pruning
on clean data pins the identified edges to their true value and leaves the
non-identifiable direct edge even, while a soft-Bayesian scorer over the same
mechanisms recovers the identified edges under \DagNoiseFlipPct\%{} outcome noise.}
\label{tab:causal-assumptions}
\begin{tabular}{l l c c c}
\toprule
Edge & Candidates & Exact post.\ & Soft post.\ & Identified \\
& & (clean) & (noisy) & \\
\midrule
Income $\to$ Stress & causal/null & 1.00 & 1.00 & \checkmark \\
Stress $\to$ Preterm & causal/null & 1.00 & 1.00 & \checkmark \\
Preterm $\to$ Bayley & causal/null & 1.00 & 1.00 & \checkmark \\
Stress $\to$ Bayley (direct) & yes/no & 0.50 & 0.50 & -- \\
\bottomrule
\end{tabular}

\end{table}

\subsection{Application: social determinants of health}\label{sec:dag-app}
The motivating use case is individual health-risk counterfactuals over social
determinants of health (SDOH). A person's record is perceived into a symbolic
SDOH state; structural factors (income, housing, food security, social support,
neighborhood) and a place-based population prior drive a mediating psychosocial
layer (perceived stress, pregnancy-related anxiety), which in turn drives a
physiological load (a neuroendocrine/cortisol and an inflammatory term) and
finally an individualized outcome risk. The chain encodes a documented
mechanism---psychosocial stress in pregnancy raising preterm-birth risk through
neuroendocrine and inflammatory pathways, with social support protective and
effects strongest in early gestation \cite{shapiro2013psp}---over a social
gradient in which structural disadvantage, not individual behavior alone, sets
the baseline \cite{marmot2018}. Interventions are the world's actions, so
counterfactual care plans are rollouts under different action sequences: the
framework compares a structural plan, a proximal stress/anxiety-reduction plan,
and their combination by resulting risk and by risk reduction per unit of effort,
and emits a recommendation. Two worked worlds ship as examples (a generic SDOH
world and a perinatal preterm-birth world); both reproduce the qualitative
literature---structural action dominates behavior-only plans, while proximal
stress care matters most early in gestation.

The scope is deliberately bounded. \openworld{} supplies the verified, auditable
counterfactual layer; it consumes a population prior (e.g.\ a geospatial
foundation-model embedding for a locale) as an exogenous input rather than
learning one. And it tests the \emph{implications and internal consistency} of a
causal assumption set---interventional and counterfactual predictions, edge-sign
and functional-form sensitivity, and the identifiability of a target effect---not
the truth of the diagram, which still requires real data and an identification
argument. What it removes is the opacity: every edge is a line of code, every
recommendation a rollout one can replay, and every assumption a hypothesis the
data is allowed to refute.
.

\section{Causal DAGs as World Specifications}\label{sec:dag}
A world model and a causal diagram are the same object viewed from two sides: a
directed acyclic graph (DAG) names variables and asserts which directly cause
which; an \openworld{} \texttt{World} carries a symbolic state ontology and a
transition that recomputes that state. \openworld{} closes the gap by accepting
a DAG as a first-class input and compiling it into a verified world, so that the
causal assumptions a domain expert is willing to commit to in a diagram become
executable, inspectable, and---crucially---falsifiable dynamics.

\subsection{A DAG as a perception modality}
DAGs enter through the same perception boundary (\S\ref{sec:language}) that
admits text, audio, and images: a \texttt{graph} modality whose
\texttt{DAGPerceptor} ingests the \texttt{dagitty}/Graphviz \texttt{.dot} text
those tools already export. A single option selects what the graph resolves
into---the three ways a world model consumes a causal structure. In
\emph{graph} mode it commits a normalized object (nodes with their
exposure/outcome/latent tags, plus edges) to state. In \emph{schema} mode it
emits the observed (non-latent) variables to extract from a record downstream:
the DAG \emph{declares what to perceive}. In \emph{world} mode it compiles the
DAG to a runnable structural causal model (SCM). The perceptor is deterministic
and code-free, so it round-trips losslessly through a portable spec like the
other declarative perceptors.

\subsection{Compiling a DAG to a verified SCM}
Compilation emits a single pure \texttt{transition(state, action)} that
recomputes each node from its parents in topological order, with an exogenous
term $u_v$ per node $v$; the default per-node form is a transparent
linear threshold over the parents, and edge weights/signs are overridable. The
emitted code passes the same gates as any synthesized dynamics
(\S\ref{sec:language})---it parses, smoke-runs sandboxed on every action, and
respects declared invariants---so the causal model is an auditable artifact, not
a black box. Interventions are simply the world's actions: a \texttt{do\_$v$}
action overrides node $v$'s structural equation and the override propagates to
its descendants within the same topological pass, realizing Pearl's do-operator
\cite{pearl2009causality}. Because the dynamics are deterministic given the
exogenous terms, a unit-level counterfactual is recovered exactly: hold a unit's
$u_v$ fixed, swap the action, and re-roll (abduction--action--prediction). The
result is a training-free interventional and counterfactual sandbox over an
\emph{assumed} causal structure.

\subsection{Testing the assumptions, not just running them}
A diagram is a set of hypotheses, and \openworld{} lets a corpus of observed
transitions adjudicate among them on the factored many-worlds store of
\S\ref{sec:active-induction}. Each uncertain edge becomes a parameter whose
candidate values are the rival assumptions---an edge's presence, or the sign of
an effect---and each node a \texttt{Mechanism} predicting a discretized
(cut-point) outcome from its parents under those parameters. Observing
transitions reweights only the small factors the firing mechanisms touch, so the
version space over causal models is pruned without ever enumerating the joint.

On a synthetic cohort of \DagCleanCohort{} units drawn from a known structure
(E64; shaped after the perinatal DAG of \S\ref{sec:dag-app}), exact pruning
collapses \DagCandidateModels{} candidate causal models to \DagSurvivorsClean{},
recovering all \DagIdentifiedEdges{} identified edges
(income~$\to$~stress, stress~$\to$~preterm, preterm~$\to$~outcome) at posterior
$1.0$ (Table~\ref{tab:causal-assumptions}). The \DagSurvivorsClean{} survivors
differ only in whether a \emph{direct} stress~$\to$~outcome edge is present:
with stress fully determining preterm in this cohort, that edge is
\emph{non-identifiable}, and the store reports it as an even posterior rather
than committing---the mediator-versus-direct ambiguity surfaced and measured,
not hidden, in the same spirit as the framework's other boundary results. Hard
version-space pruning is brittle to measurement noise: under \DagNoiseFlipPct\%{}
outcome noise the true model is contradicted and \DagHardSurvivorsNoisy{}
candidates survive. An additive soft-Bayesian scorer over the \emph{same}
mechanisms, weighting each observation by a flip-noise likelihood, instead holds
the identified edges at posterior $\geq\DagSoftMinIdentified{}$ while still
leaving the unidentifiable edge near-even---motivating a likelihood-weighted
semiring as the natural generalization of the store's update.

\begin{table}[htb]
\centering
\caption{\textbf{E64: testing causal assumptions on the many-worlds store.} Each
questioned edge of the perinatal DAG is a parameter; exact version-space pruning
on clean data pins the identified edges to their true value and leaves the
non-identifiable direct edge even, while a soft-Bayesian scorer over the same
mechanisms recovers the identified edges under \DagNoiseFlipPct\%{} outcome noise.}
\label{tab:causal-assumptions}
\begin{tabular}{l l c c c}
\toprule
Edge & Candidates & Exact post.\ & Soft post.\ & Identified \\
& & (clean) & (noisy) & \\
\midrule
Income $\to$ Stress & causal/null & 1.00 & 1.00 & \checkmark \\
Stress $\to$ Preterm & causal/null & 1.00 & 1.00 & \checkmark \\
Preterm $\to$ Bayley & causal/null & 1.00 & 1.00 & \checkmark \\
Stress $\to$ Bayley (direct) & yes/no & 0.50 & 0.50 & -- \\
\bottomrule
\end{tabular}

\end{table}

\subsection{Application: social determinants of health}\label{sec:dag-app}
The motivating use case is individual health-risk counterfactuals over social
determinants of health (SDOH). A person's record is perceived into a symbolic
SDOH state; structural factors (income, housing, food security, social support,
neighborhood) and a place-based population prior drive a mediating psychosocial
layer (perceived stress, pregnancy-related anxiety), which in turn drives a
physiological load (a neuroendocrine/cortisol and an inflammatory term) and
finally an individualized outcome risk. The chain encodes a documented
mechanism---psychosocial stress in pregnancy raising preterm-birth risk through
neuroendocrine and inflammatory pathways, with social support protective and
effects strongest in early gestation \cite{shapiro2013psp}---over a social
gradient in which structural disadvantage, not individual behavior alone, sets
the baseline \cite{marmot2018}. Interventions are the world's actions, so
counterfactual care plans are rollouts under different action sequences: the
framework compares a structural plan, a proximal stress/anxiety-reduction plan,
and their combination by resulting risk and by risk reduction per unit of effort,
and emits a recommendation. Two worked worlds ship as examples (a generic SDOH
world and a perinatal preterm-birth world); both reproduce the qualitative
literature---structural action dominates behavior-only plans, while proximal
stress care matters most early in gestation.

The scope is deliberately bounded. \openworld{} supplies the verified, auditable
counterfactual layer; it consumes a population prior (e.g.\ a geospatial
foundation-model embedding for a locale) as an exogenous input rather than
learning one. And it tests the \emph{implications and internal consistency} of a
causal assumption set---interventional and counterfactual predictions, edge-sign
and functional-form sensitivity, and the identifiability of a target effect---not
the truth of the diagram, which still requires real data and an identification
argument. What it removes is the opacity: every edge is a line of code, every
recommendation a rollout one can replay, and every assumption a hypothesis the
data is allowed to refute.
